% Modernes Abkürzungsverzeichnis für LuaLaTeX
% Nutzt glossaries statt veraltetes nomencl

% --- Glossaries-Paket (bereits in load_packages.tex geladen) ---
% \usepackage[acronym,toc,nogroupskip]{glossaries}

% --- Kompilierung mit moderne Tools ---
% LuaLaTeX + glossaries benötigt:
% 1. lualatex document.tex
% 2. makeglossaries document  (oder: makeglossaries-lite document)
% 3. lualatex document.tex
% 4. lualatex document.tex

% --- Deutsche Übersetzungen ---
\renewcommand{\glossaryname}{Glossar}
\renewcommand{\acronymname}{Abkürzungsverzeichnis}

% --- Glossar-Stile konfigurieren ---
% Moderner Stil für Abkürzungen
\setglossarystyle{long3col}              % 3-Spalten: Abkürzung | Beschreibung | Seite
% Alternativen:
% \setglossarystyle{altlong4col}         % 4-Spalten mit zusätzlicher Info
% \setglossarystyle{longtable}           % Einfache lange Tabelle
% \setglossarystyle{super3col}           % Kompakter 3-Spalten-Stil

% --- Neuer Stil mit fester Spaltenbreite ---
\newglossarystyle{fixedlong}{%
  \setglossarystyle{long}%
  \renewenvironment{theglossary}%
    {\begin{longtable}{p{0.3\textwidth}p{0.65\textwidth}}}%
    {\end{longtable}}%
}

% --- Spezifischer Stil für Acronyme ---
\setacronymstyle{long-short}             % Zeigt "Volltext (ABK)" beim ersten Mal
% Alternativen:
% \setacronymstyle{short-long}           % Zeigt "ABK (Volltext)"
% \setacronymstyle{footnote}             % Erklärung in Fußnote

% --- Layout-Anpassungen ---
\renewcommand*{\glstextformat}[1]{\textcolor{primary}{#1}}  % Farbige Links zu Glossar

% Anpassung der Spaltenbreiten
\setlength{\glsdescwidth}{0.6\textwidth}    % Beschreibungsspalte
\setlength{\glspagelistwidth}{0.1\textwidth} % Seitenzahlen-Spalte

% --- Moderne Formatierung ---
% Moderne glossaries verwenden andere Befehle für Formatierung
\renewcommand*{\glstextformat}[1]{\textcolor{primary}{#1}}  % Farbige Links zu Glossar

% --- Erweiterte Befehle für bessere Verwendung ---
% Einfacher Abkürzungsbefehl (kompatibel mit altem \abk)
\newcommand{\abk}[2]{%
  \newacronym{#1}{#1}{#2}%
}

% Für Symbole und mathematische Begriffe  
\newcommand{\defsymbol}[3][]{%
  \newglossaryentry{#2}{
    type=symbols,
    name={#2},
    description={#3},
    #1
  }%
}

% Für Fachbegriffe
\newcommand{\begriff}[3][]{%
  \newglossaryentry{#2}{
    name={#2},
    description={#3},
    #1
  }%
}

% --- Verwendung im Text (moderne Befehle) ---
% \gls{label}        - Erstes Mal: "Volltext (ABK)", danach: "ABK"
% \glspl{label}      - Plural-Form
% \Gls{label}        - Großgeschrieben am Satzanfang
% \glsfirst{label}   - Immer die Vollform
% \glsshort{label}   - Immer nur die Abkürzung
% \glslong{label}    - Immer nur die Langform


% --- Spezielle Glossar-Arten ---
% Für verschiedene Kategorien
\newglossary[slg]{symbols}{syi}{syg}{Symbolverzeichnis}
\newglossary[nlg]{notation}{not}{ntn}{Notation}

% --- Integration mit Hyperref ---
\makeatletter
\@ifpackageloaded{hyperref}{%
  % Angepasste Farben für Glossar-Links
  \renewcommand*{\glstextformat}[1]{%
    \textcolor{linkcolor}{#1}%
  }
}{}
\makeatother

% --- Sortierung und Stil ---
\GlsSetQuote{+}                          % Quote-Zeichen für Sortierung

% Deutsche Sortierung
\renewcommand*{\glsnamefont}[1]{\textbf{\textcolor{headingcolor}{#1}}}

% --- Erweiterte Darstellungsoptionen ---
% Kompakte Darstellung (weniger Zeilenabstand)
\setglossarysection{section}             % Als Section statt Chapter

% Keine Seitenzahlen (falls nicht gewünscht)
% \glsaddall                             % Alle Einträge anzeigen (auch unbenutzte)

% --- Convenience-Makros für häufige Abkürzungen ---
% Vordefinierte technische Abkürzungen
\AtBeginDocument{%
  % Nur definieren, falls noch nicht vorhanden
  \providecommand{\definecommonabbrevs}{%
    \abk{API}{Application Programming Interface}
    \abk{CPU}{Central Processing Unit}
    \abk{GPU}{Graphics Processing Unit}
    \abk{RAM}{Random Access Memory}
    \abk{SSD}{Solid State Drive}
    \abk{USB}{Universal Serial Bus}
    \abk{PDF}{Portable Document Format}
    \abk{HTML}{Hypertext Markup Language}
    \abk{CSS}{Cascading Style Sheets}
    \abk{JSON}{JavaScript Object Notation}
    \abk{XML}{Extensible Markup Language}
    \abk{HTTP}{Hypertext Transfer Protocol}
    \abk{HTTPS}{Hypertext Transfer Protocol Secure}
    \abk{URL}{Uniform Resource Locator}
    \abk{GUI}{Graphical User Interface}
    \abk{CLI}{Command Line Interface}
    \abk{AI}{Artificial Intelligence}
    \abk{ML}{Machine Learning}
    \abk{IoT}{Internet of Things}
    \abk{VPN}{Virtual Private Network}
  }
}

% Aktiviere vordefinierte Abkürzungen (optional)
% \definecommonabbrevs

% --- Mehrsprachige Unterstützung ---
\makeatletter
\@ifpackageloaded{polyglossia}{%
  % Deutsche und englische Begriffe
  \newcommand{\abkde}[2]{\newacronym[description={#2 (deutsch)}]{#1}{#1}{#2}}
  \newcommand{\abken}[2]{\newacronym[description={#2 (english)}]{#1}{#1}{#2}}
}{}
\makeatother

% --- Debug und Hilfe ---
\newcommand{\showglossarystats}{%
  \typeout{=== Glossar-Statistiken ===}
  \typeout{Definierte Abkürzungen: \the\value{glossaryentry}}
  \typeout{Verwendete Styles: \@gls@style}
  \typeout{============================}
}

% --- Finale Konfiguration ---
\makeglossaries                          % Aktiviert Glossar-Generierung

% --- Debug-Information ---
\AtBeginDocument{%
  \typeout{=== Abkürzungsverzeichnis ===}
  \typeout{System: glossaries (modern)}
  \typeout{Stil: long3col mit Farben}
  \typeout{Sprache: Deutsch + Polyglossia}
  \typeout{Kompilierung: makeglossaries benötigt}
  \typeout{Befehle: \string\abk{}{}, \string\gls{}}
  \typeout{==============================}
}

% --- Verwendungsbeispiele (auskommentiert) ---
% \abk{API}{Application Programming Interface}
% \abk{REST}{Representational State Transfer}
% 
% Im Text:
% \gls{API}          % Erstes Mal: "Application Programming Interface (API)"
% \gls{API}          % Danach: "API"
% \glspl{API}        % Plural: "APIs"
% \Gls{REST}         % Großgeschrieben: "Representational State Transfer (REST)"